\section{Conclusions}
\label{sec:conclusions}

Reservoir bathymetry and storage over the drawdown-exposed band can be reconstructed
from remote sensing alone, by stacking Sentinel-1 waterlines ordered by a remotely
supplied water level. Across nine Sicilian reservoirs the reconstruction reproduced a
December-2025 echo-sounder survey to $2.7$~m RMSE and recovered $94\%$ of the surveyed
band storage loss at that site.

Removing the in-situ gauge costs little. Rebuilding the hypsometric curve from SWOT
alone changed band storage by $0.56$~\si{Mm^3} on average, $6\%$ of band volume with a
median of $2.4\%$, and reproduced the gauge-based DEM to below $1$~m at six of the nine
reservoirs. Applied to each reservoir's full Sentinel-1 area record and scored against
officially registered volumes, the gauge-free curve outperformed the best available
independent reference at six of the nine. Because SWOT observes by swath rather than
along a nadir track, this lifts both the gauge and the ground-track constraint.

Two bounds are quantitative and should travel with any such estimate. The exposed band
holds $39$--$92\%$ of design-curve volume, so the SAR capacity-change figure is a lower
bound on total loss, deep-zone sedimentation being unobserved. Reliability scales with
the shoreline area-to-perimeter ratio, the three lowest-A/P reservoirs being the three
that reproduce the gauge-based DEM worst ($1.97$--$3.82$~m) and, at Ancipa, the one
where an independent optical reconstruction disagrees most.
